\documentclass[12pt,a4paper]{article}

\usepackage[utf8]{inputenc}
\usepackage[T1]{fontenc}
\usepackage[margin=2.5cm]{geometry}
\usepackage{amsmath,amssymb}
\usepackage{booktabs}
\usepackage{array}
\usepackage{hyperref}
\usepackage{microtype}

\hypersetup{
  colorlinks=true,
  linkcolor=blue!60!black,
  urlcolor=blue!60!black,
}

\title{\textbf{Datathon 2026 --- Ensemble Forecasting of Revenue and COGS}}
\author{}
\date{May 2026}

\begin{document}

\section{Revenue \& COGS Forecast 2023--2024}
\begin{center}
\centering
\captionof{table}{Forecast for Revenue and COGS in 2023--2024.}
\label{tab:forecast-summary}
\begin{tabular}{ll}
\toprule
\textbf{Metric} & \textbf{Value} \\
\midrule
Total Revenue & 2.322B VND \\
Daily Average & 4,238,128 VND \\
Total COGS & 2.109B VND \\
Forecast Accuracy ($R^2$) & 0.7778 (77.8\% variance) \\
Typical Daily Error (MAE) & $\pm$523K VND ($\pm$12\%) \\
\bottomrule
\end{tabular}
\end{center}

\subsection*{What Drives Revenue? (5 Key Patterns)}

\begin{center}
\centering
\captionof{table}{Main revenue drivers identified from the forecasting model.}
\label{tab:revenue-drivers}
\small
\begin{tabular}{p{4cm}p{4.2cm}p{5cm}}
\toprule
\textbf{Pattern} & \textbf{Observed evidence} & \textbf{Business interpretation} \\
\midrule
Annual Seasonality & Q3 +18\%, Q1/Q4 --8\% & Summer peaks versus winter lows \\
Weekly Cycles & Wed--Thu +9\% above average & Work-schedule effect \\
Month-End Spikes & +13\% boost at month-end & Payroll-driven purchasing \\
Long-Term Growth & 3--5\% annual expansion & Growth continues into 2023--2024 \\
Daily Randomness & Promotions, shocks, and events & Additional $\pm$12\% error margin \\
\bottomrule
\end{tabular}
\end{center}

\subsection*{Method: Theta Decomposition (17.4\%) + Hybrid ARIMA (82.6\%)}

Both models capture \textbf{long-term growth explicitly}---why classical time-series beats tree/neural methods that cap at training range \cite{Hyndman2021}. The ensemble approach is motivated by ensemble validation theory \cite{Makridakis2020}. The reliability evidence, validation protocol, data setting, and reproducibility controls are reported in the Appendix.

\subsection*{Business Takeaways}

\begin{center}
\centering
\captionof{table}{Business interpretation of the forecast results.}
\label{tab:business-takeaways}
\small
\begin{tabular}{p{3.5cm}p{10cm}}
\toprule
\textbf{Area} & \textbf{Takeaway} \\
\midrule
Revenue & \textbf{Stable} at 4.24M/day; growth continues modestly \\
COGS ratio & \textbf{Healthy} at 90.8\%; operational efficiency maintained \\
Seasonality & \textbf{Predictable} ($\pm$18\% annual, $\pm$9\% weekly, +13\% month-end) \\
Error margin & \textbf{Manageable} ($\pm$523K/day; plan 12\% buffer for daily operations) \\
Model & \textbf{Proven} on 3-year holdout; not overfit \\
\bottomrule
\end{tabular}
\end{center}

\label{sec:part3}

// Appendix
\section{Forecasting Model: Robustness and Technical Evidence}
\label{sec:forecasting-model-technical-evidence}

This section consolidates the supporting technical evidence for the Revenue and COGS forecasting model. It reports the historical validation results, experimental protocol, preprocessing design, cross-validation methodology, leakage-control rules, explainability checks, mathematical formulation, hyperparameter settings, model-selection rationale, and day-of-week factors used to justify the reliability of the forecast results presented in the main analysis.

\subsection{Historical Validation and Reliability}
\label{sec:historical-validation}

\begin{center}
\centering
\captionof{table}{Historical validation performance across three holdout years.}
\label{tab:historical-validation}
\small
\begin{tabular}{lrr}
\toprule
\textbf{Year Tested} & \textbf{Forecast Error (MAE)} & \textbf{Variance Captured ($R^2$)} \\
\midrule
2020 & 520K VND & 0.784 \\
2021 & 508K VND & 0.770 \\
2022 & 540K VND & 0.779 \\
\bottomrule
\end{tabular}
\end{center}

\noindent\textbf{Why it works}: Consistent $R^2 \approx 0.78$ across three years proves the model captures repeatable seasonal patterns and long-term growth trend---not overfitting. Expanding-window CV, where the model is trained strictly before each test year, prevents test leakage.

\subsection{Experimental Protocol and Reproducibility}
\label{sec:protocol-reproducibility}

\begin{center}
\centering
\captionof{table}{Summary of validation, data, and reproducibility settings.}
\label{tab:protocol-reproducibility}
\small
\begin{tabular}{p{3.5cm}p{10cm}}
\toprule
\textbf{Item} & \textbf{Setting} \\
\midrule
Validation & Expanding-window CV, trained strictly before each test year \\
Data & \texttt{sales.csv} only, 2012--2022, no external data \\
Reproducibility & Fixed seed: 42; full code in \texttt{scripts/forecast.py} \\
\bottomrule
\end{tabular}
\end{center}

\subsection{Input Preprocessing and Feature Engineering}
\label{sec:preprocessing-feature-engineering}

\subsubsection{Data Handling}

Daily sales data were sorted chronologically. Missing values within training windows were linearly interpolated and then forward/backward-filled. Interpolation was applied independently per fold to prevent future-data leakage. All targets were clipped to $\geq 0$ before model input.

\subsubsection{Features Used}

\begin{center}
\centering
\captionof{table}{Feature engineering design for the forecasting models.}
\label{tab:features-used}
\small
\begin{tabular}{llll}
\toprule
\textbf{Model} & \textbf{Feature} & \textbf{Encoding} & \textbf{Purpose} \\
\midrule
Theta & Revenue/COGS & $\log(1+y)$ & Variance stabilization \\
Theta & Day-of-week & Multiplicative factor & Weekly rhythm \\
ARIMA & Month, day & Calendar profile & Intra-year baseline \\
ARIMA & Year index & Geometric growth & Trend extrapolation \\
ARIMA & 3-year residuals & ARIMA fit & Autocorrelation \\
\bottomrule
\end{tabular}
\end{center}

The Theta method applies log transformation for variance stabilization and incorporates day-of-week factors \cite{Assimakopoulos2000}, while ARIMA focuses on residual patterns over a 3-year window to capture short-term autocorrelation.

\subsubsection{Cyclic Encoding}

\[
\mathrm{doy\_sin} = \sin\left(\frac{2\pi \cdot \mathrm{doy}}{365.25}\right),
\qquad
\mathrm{doy\_cos} = \cos\left(\frac{2\pi \cdot \mathrm{doy}}{365.25}\right)
\]

\noindent\textbf{Rationale}: Raw integer features, where $\mathrm{doy} \in [1,365]$, create false discontinuity at the year boundary. Cyclic sine/cosine encoding preserves ordinal distance.

\subsection{Cross-Validation Methodology}
\label{sec:cross-validation-methodology}

\begin{center}
\centering
\captionof{table}{Expanding-window cross-validation design.}
\label{tab:expanding-window-cv}
\small
\begin{tabular}{lll}
\toprule
\textbf{Fold} & \textbf{Training Period} & \textbf{Validation Period} \\
\midrule
2020 & 2012-07-04 to 2019-12-31 & 2020 \\
2021 & 2012-07-04 to 2020-12-31 & 2021 \\
2022 & 2012-07-04 to 2021-12-31 & 2022 \\
\bottomrule
\end{tabular}
\end{center}

Expanding-window validation was selected instead of rolling-window validation because rolling windows discard historical information and may weaken the estimation of long-run growth. In contrast, expanding-window validation preserves all available past observations while ensuring that each validation year is forecast using only earlier data \cite{Bergmeir2014}.

\subsection{Data Leakage Control}
\label{sec:data-leakage-control}

\begin{center}
\centering
\captionof{table}{Controls applied to prevent data leakage during validation and forecasting.}
\label{tab:data-leakage-control}
\small
\begin{tabular}{p{5.2cm}p{8.3cm}}
\toprule
\textbf{Control Measure} & \textbf{Implementation} \\
\midrule
No future Revenue/COGS in training & Strict filter: Date $<$ fold year \\
No global statistics from test data & Scales and day-of-week profiles computed per fold only \\
Day-of-week factors recomputed per fold & Each fold uses its training data only \\
Interpolation within fold only & Missing values filled from the training window only \\
Ensemble weights from out-of-fold validation & Weights optimized on fold validation results \\
Sample submission used for dates only & Revenue/COGS columns ignored \\
\bottomrule
\end{tabular}
\end{center}

\subsection{Feature Importance and Explainability}
\label{sec:feature-importance-explainability}

Absolute Pearson correlation was calculated using concatenated Revenue and COGS values to provide a simple interpretability check for the forecasting features.

\begin{center}
\centering
\captionof{table}{Feature importance estimated by absolute Pearson correlation.}
\label{tab:feature-importance}
\small
\begin{tabular}{lrl}
\toprule
\textbf{Feature} & \textbf{Corr. (\%)} & \textbf{Interpretation} \\
\midrule
doy\_cos & 16.86 & Cyclic annual phase, seasonal high \\
is\_month\_end & 13.17 & End-of-month purchasing spike \\
days\_since & 10.91 & Long-term growth trend \\
doy\_sin & 10.68 & Cyclic annual phase, Jan--Jul movement \\
year & 10.31 & Year-over-year business growth \\
dom & 9.81 & Day-of-month position \\
month & 6.87 & Broad seasonal periods \\
week & 6.36 & Within-year seasonal progression \\
\bottomrule
\end{tabular}
\end{center}

The most important features correspond to interpretable business patterns: yearly seasonality, month-end purchasing behavior, long-term growth, and calendar-driven demand variation.

\subsection{Mathematical Formulas}
\label{sec:mathematical-formulas}

\subsubsection{Competition Metrics}

The three evaluation metrics used throughout this analysis are computed as follows \cite{Pedregosa2011}:

\[
\mathrm{MAE} = \frac{1}{n} \sum_{i=1}^{n} |F_i - A_i|,
\qquad
\mathrm{RMSE} = \sqrt{\frac{1}{n} \sum_{i=1}^{n} (F_i - A_i)^2},
\qquad
R^2 = 1 - \frac{\sum_{i=1}^{n}(A_i - F_i)^2}{\sum_{i=1}^{n}(A_i - \bar{A})^2}
\]

where $n = 1,096$ because Revenue and COGS are concatenated for evaluation. Numerical implementations are provided by scikit-learn.

\subsubsection{Theta Method}

The Theta method decomposes $y_t$ into two theta lines, a decomposition approach that won the M3/M4 competitions \cite{Assimakopoulos2000}:

\[
\tilde{y}_0(t) = 2\bar{y} - y_t,
\qquad
\tilde{y}_2(t) = y_t
\]

The final forecast is computed as:

\[
F(h) = \frac{1}{2}\mathrm{SES}(\tilde{y}_0, h) + \frac{1}{2}(a + b(T+h))
\]

The method was applied on $\log(1+y)$ with a learned day-of-week multiplier.

\subsubsection{Hybrid ARIMA}

The Hybrid ARIMA model combines a seasonal baseline with residual ARIMA correction.

\[
\hat{y}(t) = b \cdot g^{\Delta_{\mathrm{yr}}} \cdot s(\mathrm{month}, \mathrm{day})
\]

\[
e_t = c + \sum_{i=1}^{p} \varphi_i e_{t-i} + \sum_{j=1}^{q} \theta_j \varepsilon_{t-j} + \varepsilon_t
\]

\[
F(h) = \max(0, \hat{y}(h) + \mathrm{ARIMA\_forecast}(h))
\]

\subsection{Hyperparameters and Tuning}
\label{sec:hyperparameters-tuning}

\begin{center}
\centering
\captionof{table}{Hyperparameter configuration and tuning method.}
\label{tab:hyperparameters}
\small
\begin{tabular}{lll}
\toprule
\textbf{Parameter} & \textbf{Value} & \textbf{Tuning Method} \\
\midrule
ARIMA Revenue $(p,d,q)$ & $(2,0,3)$ & Grid search on CV \\
ARIMA COGS $(p,d,q)$ & $(2,0,3)$ & Grid search on CV \\
ARIMA residual window & 3 years & CV validation \\
Holt-Winters trend & Additive & Fixed \\
Holt-Winters seasonality & Multiplicative & Fixed \\
HW seasonal periods & 7 & Fixed, weekly \\
Theta period & 365 & Fixed, annual \\
Theta log transform & True & Fixed \\
SCALE\_REVENUE & 1.184 & CV calibration \\
SCALE\_COGS & 1.191 & CV calibration \\
DOW\_STRENGTH & 0.40 & Manual tuning \\
REVENUE\_BIAS & 25,000 VND/day & CV calibration \\
COGS\_BIAS & 112,500 VND/day & CV calibration \\
MAX\_COGS\_RATIO & 1.05 & Competition constraint \\
Optimizer & Nelder-Mead & Fixed \\
Optimizer restarts & 12 & Fixed \\
Random seed & 42 & Fixed \\
\bottomrule
\end{tabular}
\end{center}

\subsection{Model Comparison and Selection}
\label{sec:model-comparison-selection}

\begin{center}
\centering
\captionof{table}{Candidate model comparison and final ensemble selection.}
\label{tab:model-comparison}
\small
\begin{tabular}{llrr}
\toprule
\textbf{Model} & \textbf{Type} & \textbf{CV MAE} & \textbf{Decision} \\
\midrule
\textbf{Hybrid ARIMA} & Classical TS & $\sim$540K & Selected, 83\% \\
\textbf{Theta} & Classical TS & $\sim$540K & Selected, 17\% \\
Holt-Winters & Classical TS & $\sim$540K & 0\% weight \\
Prophet & Bayesian TS & $\sim$540K & 0\% weight \\
LightGBM & Tree Ensemble & $\sim$631K & Rejected \\
Chronos-T5 & Foundation Model & $\sim$1,370K & Rejected \\
\bottomrule
\end{tabular}
\end{center}

Tree-based and neural models tend to cap predictions within the training range. Classical time-series methods with explicit trend components, including geometric growth and linear extrapolation, provide more reliable extrapolation beyond the 2012--2022 training distribution.

\subsection{Day-of-Week Factors}
\label{sec:day-of-week-factors}

\begin{center}
\centering
\captionof{table}{Estimated day-of-week factors for Revenue and COGS.}
\label{tab:dow-factors}
\small
\begin{tabular}{lrr}
\toprule
\textbf{Day} & \textbf{Revenue Factor} & \textbf{COGS Factor} \\
\midrule
Monday & 1.015 & 1.013 \\
Tuesday & 1.033 & 1.032 \\
Wednesday & 1.091 & 1.091 \\
Thursday & 1.038 & 1.038 \\
Friday & 0.945 & 0.945 \\
Saturday & 0.918 & 0.919 \\
Sunday & 0.959 & 0.960 \\
\bottomrule
\end{tabular}
\end{center}

Mid-week demand, especially Wednesday and Thursday, consistently outperforms weekends by approximately 8--9\%. This pattern suggests B2B or work-hour-adjacent purchasing behavior and supports the inclusion of day-of-week factors in the forecasting model.

\newpage

\begin{thebibliography}{99}

\bibitem{Hyndman2021}
Hyndman, R. J. and Athanasopoulos, G.
\newblock \textit{Forecasting: Principles and Practice} (3rd ed.).
\newblock OTexts.
\newblock Available at: \url{https://otexts.com/fpp3/}

\bibitem{Assimakopoulos2000}
Assimakopoulos, V. and Nikolopoulos, K.
\newblock The Theta model: a decomposition approach to forecasting.
\newblock \textit{International Journal of Forecasting}, 16(4):521--530.
\newblock doi: 10.1016/S0169-2070(00)00066-2

\bibitem{Bergmeir2014}
Bergmeir, C., Ben\'itez, M., and Artieda, J.
\newblock On the use of cross-validation for time series forecasting evaluation.
\newblock \textit{Information Sciences}, 191:192--213.
\newblock doi: 10.1016/j.ins.2013.07.019

\bibitem{Pedregosa2011}
Pedregosa, F., Varoquaux, G., Gramfort, A., et al.
\newblock scikit-learn: Machine learning in Python.
\newblock \textit{Journal of Machine Learning Research}, 12:2825--2830.
\newblock Available at: \url{https://jmlr.org/papers/v12/pedregosa11a.html}

\bibitem{Harris2020}
Harris, C. R., Millman, K. J., van der Walt, S. J., et al.
\newblock Array programming with NumPy.
\newblock \textit{Nature}, 585:357--362.
\newblock doi: 10.1038/s41586-020-2649-2

\bibitem{Makridakis2020}
Makridakis, S., Spiliotis, E., and Assimakopoulos, V.
\newblock The M4 Competition: 100,000 time series and 61 forecasting methods.
\newblock \textit{International Journal of Forecasting}, 36(1):54--74.
\newblock doi: 10.1016/j.ijforecast.2019.12.011

\end{thebibliography}

\end{document}
